% Options for packages loaded elsewhere
\PassOptionsToPackage{unicode}{hyperref}
\PassOptionsToPackage{hyphens}{url}
\PassOptionsToPackage{dvipsnames,svgnames,x11names}{xcolor}
%
\documentclass[
  11pt,
]{article}
\usepackage{amsmath,amssymb}
\usepackage{iftex}
\ifPDFTeX
  \usepackage[T1]{fontenc}
  \usepackage[utf8]{inputenc}
  \usepackage{textcomp} % provide euro and other symbols
\else % if luatex or xetex
  \usepackage{unicode-math} % this also loads fontspec
  \defaultfontfeatures{Scale=MatchLowercase}
  \defaultfontfeatures[\rmfamily]{Ligatures=TeX,Scale=1}
\fi
\usepackage{lmodern}
\ifPDFTeX\else
  % xetex/luatex font selection
  \setmainfont[]{DejaVu Serif}
  \setmonofont[]{DejaVu Sans Mono}
\fi
% Use upquote if available, for straight quotes in verbatim environments
\IfFileExists{upquote.sty}{\usepackage{upquote}}{}
\IfFileExists{microtype.sty}{% use microtype if available
  \usepackage[]{microtype}
  \UseMicrotypeSet[protrusion]{basicmath} % disable protrusion for tt fonts
}{}
\makeatletter
\@ifundefined{KOMAClassName}{% if non-KOMA class
  \IfFileExists{parskip.sty}{%
    \usepackage{parskip}
  }{% else
    \setlength{\parindent}{0pt}
    \setlength{\parskip}{6pt plus 2pt minus 1pt}}
}{% if KOMA class
  \KOMAoptions{parskip=half}}
\makeatother
\usepackage{xcolor}
\usepackage[margin=1in]{geometry}
\usepackage{color}
\usepackage{fancyvrb}
\newcommand{\VerbBar}{|}
\newcommand{\VERB}{\Verb[commandchars=\\\{\}]}
\DefineVerbatimEnvironment{Highlighting}{Verbatim}{commandchars=\\\{\}}
% Add ',fontsize=\small' for more characters per line
\newenvironment{Shaded}{}{}
\newcommand{\AlertTok}[1]{\textcolor[rgb]{1.00,0.00,0.00}{\textbf{#1}}}
\newcommand{\AnnotationTok}[1]{\textcolor[rgb]{0.38,0.63,0.69}{\textbf{\textit{#1}}}}
\newcommand{\AttributeTok}[1]{\textcolor[rgb]{0.49,0.56,0.16}{#1}}
\newcommand{\BaseNTok}[1]{\textcolor[rgb]{0.25,0.63,0.44}{#1}}
\newcommand{\BuiltInTok}[1]{\textcolor[rgb]{0.00,0.50,0.00}{#1}}
\newcommand{\CharTok}[1]{\textcolor[rgb]{0.25,0.44,0.63}{#1}}
\newcommand{\CommentTok}[1]{\textcolor[rgb]{0.38,0.63,0.69}{\textit{#1}}}
\newcommand{\CommentVarTok}[1]{\textcolor[rgb]{0.38,0.63,0.69}{\textbf{\textit{#1}}}}
\newcommand{\ConstantTok}[1]{\textcolor[rgb]{0.53,0.00,0.00}{#1}}
\newcommand{\ControlFlowTok}[1]{\textcolor[rgb]{0.00,0.44,0.13}{\textbf{#1}}}
\newcommand{\DataTypeTok}[1]{\textcolor[rgb]{0.56,0.13,0.00}{#1}}
\newcommand{\DecValTok}[1]{\textcolor[rgb]{0.25,0.63,0.44}{#1}}
\newcommand{\DocumentationTok}[1]{\textcolor[rgb]{0.73,0.13,0.13}{\textit{#1}}}
\newcommand{\ErrorTok}[1]{\textcolor[rgb]{1.00,0.00,0.00}{\textbf{#1}}}
\newcommand{\ExtensionTok}[1]{#1}
\newcommand{\FloatTok}[1]{\textcolor[rgb]{0.25,0.63,0.44}{#1}}
\newcommand{\FunctionTok}[1]{\textcolor[rgb]{0.02,0.16,0.49}{#1}}
\newcommand{\ImportTok}[1]{\textcolor[rgb]{0.00,0.50,0.00}{\textbf{#1}}}
\newcommand{\InformationTok}[1]{\textcolor[rgb]{0.38,0.63,0.69}{\textbf{\textit{#1}}}}
\newcommand{\KeywordTok}[1]{\textcolor[rgb]{0.00,0.44,0.13}{\textbf{#1}}}
\newcommand{\NormalTok}[1]{#1}
\newcommand{\OperatorTok}[1]{\textcolor[rgb]{0.40,0.40,0.40}{#1}}
\newcommand{\OtherTok}[1]{\textcolor[rgb]{0.00,0.44,0.13}{#1}}
\newcommand{\PreprocessorTok}[1]{\textcolor[rgb]{0.74,0.48,0.00}{#1}}
\newcommand{\RegionMarkerTok}[1]{#1}
\newcommand{\SpecialCharTok}[1]{\textcolor[rgb]{0.25,0.44,0.63}{#1}}
\newcommand{\SpecialStringTok}[1]{\textcolor[rgb]{0.73,0.40,0.53}{#1}}
\newcommand{\StringTok}[1]{\textcolor[rgb]{0.25,0.44,0.63}{#1}}
\newcommand{\VariableTok}[1]{\textcolor[rgb]{0.10,0.09,0.49}{#1}}
\newcommand{\VerbatimStringTok}[1]{\textcolor[rgb]{0.25,0.44,0.63}{#1}}
\newcommand{\WarningTok}[1]{\textcolor[rgb]{0.38,0.63,0.69}{\textbf{\textit{#1}}}}
\usepackage{longtable,booktabs,array}
\usepackage{calc} % for calculating minipage widths
% Correct order of tables after \paragraph or \subparagraph
\usepackage{etoolbox}
\makeatletter
\patchcmd\longtable{\par}{\if@noskipsec\mbox{}\fi\par}{}{}
\makeatother
% Allow footnotes in longtable head/foot
\IfFileExists{footnotehyper.sty}{\usepackage{footnotehyper}}{\usepackage{footnote}}
\makesavenoteenv{longtable}
\setlength{\emergencystretch}{3em} % prevent overfull lines
\providecommand{\tightlist}{%
  \setlength{\itemsep}{0pt}\setlength{\parskip}{0pt}}
\setcounter{secnumdepth}{-\maxdimen} % remove section numbering
\ifLuaTeX
  \usepackage{selnolig}  % disable illegal ligatures
\fi
\IfFileExists{bookmark.sty}{\usepackage{bookmark}}{\usepackage{hyperref}}
\IfFileExists{xurl.sty}{\usepackage{xurl}}{} % add URL line breaks if available
\urlstyle{same}
\hypersetup{
  pdftitle={Corda (2021) `The secret of planets' perihelion between Newton and Einstein' --- math + physical-intuition check for the Relativistic Newtonian Theory foundations},
  pdfauthor={Trey Morris with Claude Opus 4.7},
  colorlinks=true,
  linkcolor={blue},
  filecolor={Maroon},
  citecolor={Blue},
  urlcolor={blue},
  pdfcreator={LaTeX via pandoc}}

\title{Corda (2021) `The secret of planets' perihelion between Newton
and Einstein' --- math + physical-intuition check for the
\emph{Relativistic Newtonian Theory} foundations}
\author{Trey Morris with Claude Opus 4.7}
\date{2026-05-31}

\begin{document}
\maketitle

\hypertarget{author-review-report-corda-2021-the-secret-of-planets-perihelion-between-newton-and-einstein-math-physical-intuition-check}{%
\section{Author-review report: Corda (2021) ``The secret of planets'
perihelion between Newton and Einstein'' --- math + physical-intuition
check}\label{author-review-report-corda-2021-the-secret-of-planets-perihelion-between-newton-and-einstein-math-physical-intuition-check}}

\textbf{To:} Tepper Gill (DRQM I + \emph{Relativistic Newtonian Theory}
senior author) \textbf{From:} Trey Morris \textbf{Date:} 2026-05-30
\textbf{Re:} Issue
\href{https://github.com/temoTxt/PyPhysics/issues/81}{\#81}, basis for
Tepper's Mercury perihelion derivation \textbf{Source paper under
review:} Christian Corda, \emph{Phys. Dark Univ.} \textbf{32} (2021)
100834,
\href{https://doi.org/10.1016/j.dark.2021.100834}{doi:10.1016/j.dark.2021.100834}.
PDF in repo at
\href{../External_Papers/Secret_Corda.pdf}{\texttt{Roadmapping/External\_Papers/Secret\_Corda.pdf}};
marker-pdf conversion at
\href{../Converted_Markdown/Secret_Corda/Secret_Corda.md}{\texttt{Roadmapping/Converted\_Markdown/Secret\_Corda/Secret\_Corda.md}}.

\begin{center}\rule{0.5\linewidth}{0.5pt}\end{center}

\hypertarget{cover-note-for-author-review-not-for-journal-submission}{%
\subsection{Cover note (for author review, not for journal
submission)}\label{cover-note-for-author-review-not-for-journal-submission}}

You indicated that the perihelion derivation in your \emph{Relativistic
Newtonian Theory} paper takes Corda's framework as its starting point.
This report checks Corda's math and physical intuition end-to-end so you
can adjudicate whether the framework's intended Mercury prediction is
the Corda-style \texttt{πm/M\ +\ relativistic\ correction} template (the
chain analysed in PR
\href{https://github.com/temoTxt/PyPhysics/pull/83}{\#83} giving
\texttt{+37.79\textquotesingle{}\textquotesingle{}/century}) or the
proper apsidal-angle calculation of the modified force law (PR \#83 doc
05 giving \texttt{−7.17\textquotesingle{}\textquotesingle{}/century}).

Per CLAUDE.md author-review convention, the findings below are surfaced
for your verdict --- not announced as corrections. Every algebraic claim
is verified symbolically via Wolfram MCP; every interpretive claim is
flagged as substantive AI and tagged with author-review markers
\texttt{{[}for\ Tepper:{]}}.

Per Crocco §5 compliance, this is substantive AI. The prompt-of-record
is the user instruction sequence on issue \#81 and the user's 2026-05-30
ask ``write the complete analysis in a markdown file for further review
by Tepper.'' No prose generation prompt beyond the user's instructions
was used.

\begin{center}\rule{0.5\linewidth}{0.5pt}\end{center}

\hypertarget{cordas-paper-structure-8-pages-77-numbered-equations}{%
\subsection{§1. Corda's paper --- structure (8 pages, 77 numbered
equations)}\label{cordas-paper-structure-8-pages-77-numbered-equations}}

\begin{longtable}[]{@{}
  >{\raggedright\arraybackslash}p{(\columnwidth - 6\tabcolsep) * \real{0.2500}}
  >{\raggedright\arraybackslash}p{(\columnwidth - 6\tabcolsep) * \real{0.2500}}
  >{\raggedright\arraybackslash}p{(\columnwidth - 6\tabcolsep) * \real{0.2500}}
  >{\raggedright\arraybackslash}p{(\columnwidth - 6\tabcolsep) * \real{0.2500}}@{}}
\toprule\noalign{}
\begin{minipage}[b]{\linewidth}\raggedright
§
\end{minipage} & \begin{minipage}[b]{\linewidth}\raggedright
Topic
\end{minipage} & \begin{minipage}[b]{\linewidth}\raggedright
Headline equation
\end{minipage} & \begin{minipage}[b]{\linewidth}\raggedright
Result
\end{minipage} \\
\midrule\noalign{}
\endhead
\bottomrule\noalign{}
\endlastfoot
§1 & Introduction, GR perihelion & Eq. (1):
\texttt{Δφ\_GR\ =\ 24π³\ a²\ /\ {[}T₀²\ c²\ (1−e²){]}\ =\ 3π\ r\_g\ /\ {[}a(1−e²){]}}
& \texttt{42.99\textquotesingle{}\textquotesingle{}/c} for Mercury \\
§2 & ``Circular orbit'' Newtonian back-reaction & Eq. (20):
\texttt{Δφ\ =\ πm/M} &
\texttt{44.66\textquotesingle{}\textquotesingle{}/c} (Approach 1) \\
§3 & Harmonic-oscillator / apsidal-angle & Eq. (31):
\texttt{φ₀\ =\ 2π{[}3\ +\ F\textquotesingle{}(r)r/F(r){]}\^{}(−1/2)} &
same \texttt{44.66\textquotesingle{}\textquotesingle{}/c} (Approach
2) \\
§4 & Reduced-mass Kepler & Eq. (57):
\texttt{T\ =\ T₀\ (1+m/M)\^{}(−1/2)} & same
\texttt{44.66\textquotesingle{}\textquotesingle{}/c} (Approach 3) \\
§5 & Breakdown of \texttt{πm/M} for Venus / Earth & (numerical only) &
\texttt{30×} and \texttt{51×} too large \\
§6 & Gravitational time dilation & Eq. (67):
\texttt{Δφ\_F\ =\ (5π/2)\ r\_g/r₀} &
\texttt{34.31\textquotesingle{}\textquotesingle{}/c} (TD-only) \\
§7 & Rotational (Langevin) time dilation & Eq. (72):
\texttt{τ\ ≈\ (1\ −\ ½\ r²ω²/c²)\ t} & (intermediate) \\
§8 & Combined gravitational + rotational TD & Eq. (77):
\texttt{Δφ\_F\ =\ 3π\ r\_g/r₀} &
\texttt{41.17\textquotesingle{}\textquotesingle{}/c} (matches e=0 GR) \\
§9 & Conclusion: ``Newtonian + corrections matches GR'' & --- & --- \\
\end{longtable}

\textbf{The paper's two thesis claims (§9 verbatim):}

\begin{quote}
\begin{enumerate}
\def\labelenumi{(\roman{enumi})}
\item
  It is not correct that Newtonian theory cannot predict the anomalous
  rate of precession of the perihelion of planets' orbit. The real
  problem is instead that a pure Newtonian prediction is too large.
\item
  Perihelion's precession can be achieved with the same precision of
  general relativity by extending Newtonian gravity through the
  inclusion of gravitational and rotational time dilation effects.
\end{enumerate}
\end{quote}

The math + physical-intuition check below addresses each thesis
separately.

\begin{center}\rule{0.5\linewidth}{0.5pt}\end{center}

\hypertarget{every-load-bearing-equation-surfaced-and-verified}{%
\subsection{§2. Every load-bearing equation surfaced and
verified}\label{every-load-bearing-equation-surfaced-and-verified}}

Each subsection below states Corda's equation verbatim (translated from
the marker-pdf conversion), notes the Wolfram-MCP-verified algebra, and
flags the physical-intuition status.

\hypertarget{eqs.-35-circular-orbit-newton-baseline-no-back-reaction}{%
\subsubsection{§2.1 Eqs. (3)--(5) --- circular-orbit Newton baseline (no
back-reaction)}\label{eqs.-35-circular-orbit-newton-baseline-no-back-reaction}}

\[
\frac{GMm}{r_0^2} = \frac{m v_0^2}{r_0}     \qquad (3)
\] \[
v_0 = \sqrt{GM/r_0}                          \qquad (4)
\] \[
T_0 = \frac{2\pi r_0}{v_0} = \frac{2\pi r_0^{3/2}}{\sqrt{GM}}  \qquad (5)
\] Standard textbook (Newton + centripetal). {[}OK{]} Algebraically
correct. Physical content: Mercury orbits a stationary Sun. Standard.

\hypertarget{eq.-67-newtonian-angular-velocity-and-unperturbed-angle}{%
\subsubsection{§2.2 Eq. (6)--(7) --- Newtonian angular velocity and
unperturbed
angle}\label{eq.-67-newtonian-angular-velocity-and-unperturbed-angle}}

\[
\omega_0 = 2\pi / T_0                         \qquad (6)
\] \[
\varphi_0 = \omega_0 T_0 = 2\pi               \qquad (7)
\] Definitionally correct. {[}OK{]}

\hypertarget{eqs.-813-the-substitution-that-drives-cordas-2-result}{%
\subsubsection{\texorpdfstring{§2.3 Eqs. (8)--(13) --- the
\textbf{substitution} that drives Corda's §2
result}{§2.3 Eqs. (8)--(13) --- the substitution that drives Corda's §2 result}}\label{eqs.-813-the-substitution-that-drives-cordas-2-result}}

Corda's \emph{core substitution}: replace \texttt{M} in Eq. (3) with
\texttt{M\ +\ m} because the Sun-centred observer sees Mercury's
reaction force on the Sun too.

\[
\frac{G(M+m)m}{r_0^2} = \frac{m v^2}{r_0}     \qquad (8)
\] He justifies this by writing the Newton's-third-law pair (Eqs.
9--13):

\[
\vec F_G = \frac{GMm}{r_0^2}\hat u_r           \qquad (9)
\] \[
M a_s\hat u_r = \frac{GMm}{r_0^2}\hat u_r \;\Rightarrow\; a_s = \frac{Gm}{r_0^2}  \qquad (10)
\] \[
m a_m\hat u_r = -\frac{GMm}{r_0^2}\hat u_r \;\Rightarrow\; a_m = -\frac{GM}{r_0^2}  \qquad (11)
\] \[
a\hat u_r \equiv a_m - a_s = -\frac{G(M+m)}{r_0^2}\hat u_r  \qquad (12)
\] \[
F\hat u_r = -\frac{G(M+m)m}{r_0^2}\hat u_r     \qquad (13)
\] \textbf{Algebra {[}OK{]}} --- this is the standard two-body
relative-acceleration result. The same \texttt{G(M+m)} shows up in
textbook reduced-mass treatments (Goldstein §3.1, MTW §25). What
\texttt{(M+m)} captures is the kinematic fact that the \emph{relative}
acceleration of Mercury w.r.t. the Sun is bigger than \texttt{GM/r²} by
the factor \texttt{(M+m)/M\ =\ 1\ +\ m/M}.

\textbf{Physical intuition status {[}for Tepper:{]}} what is being
computed here is the \textbf{relative} acceleration in the
centre-of-mass frame, not a ``Sun-fixed'' force on Mercury. The standard
reduced-mass two-body decomposition replaces the two-body problem with a
one-body problem for a reduced-mass particle \texttt{μ\ =\ Mm/(M+m)}
orbiting a fixed centre at \texttt{M\ +\ m}. Corda's Eq. (12) is
correct, but the interpretation in §2 (that this constitutes a
``perturbation'' to Mercury's force) muddles the standard story.

\hypertarget{eqs.-1417-perturbed-velocity-and-period}{%
\subsubsection{§2.4 Eqs. (14)--(17) --- perturbed velocity and
period}\label{eqs.-1417-perturbed-velocity-and-period}}

\[
v = \sqrt{G(M+m)/r_0}                          \qquad (14)
\] \[
T = \frac{2\pi r_0}{v} = \frac{2\pi r_0^{3/2}}{\sqrt{G(M+m)}}  \qquad (15)
\] \[
(M+m)^{-1/2} = M^{-1/2}\,(1 + m/M)^{-1/2}      \qquad (16)
\] \[
T = T_0\,(1 + m/M)^{-1/2}                      \qquad (17)
\]

\textbf{Wolfram MCP check:}

\begin{Shaded}
\begin{Highlighting}[]
\FunctionTok{ClearAll}\OperatorTok{[}\NormalTok{GG}\OperatorTok{,}\NormalTok{ Msun}\OperatorTok{,}\NormalTok{ mMerc}\OperatorTok{,}\NormalTok{ r0}\OperatorTok{,} \FunctionTok{T}\OperatorTok{,}\NormalTok{ T0}\OperatorTok{]}\NormalTok{;}
\FunctionTok{v}   \ExtensionTok{=} \FunctionTok{Sqrt}\OperatorTok{[}\NormalTok{GG (Msun }\SpecialCharTok{+}\NormalTok{ mMerc)}\SpecialCharTok{/}\NormalTok{r0}\OperatorTok{]}\NormalTok{;}
\FunctionTok{T}   \ExtensionTok{=} \DecValTok{2} \FunctionTok{Pi}\NormalTok{ r0 }\SpecialCharTok{/} \FunctionTok{v}\NormalTok{;}
\NormalTok{T0  }\ExtensionTok{=} \DecValTok{2} \FunctionTok{Pi}\NormalTok{ r0}\SpecialCharTok{\^{}}\NormalTok{(}\DecValTok{3}\SpecialCharTok{/}\DecValTok{2}\NormalTok{) }\SpecialCharTok{/} \FunctionTok{Sqrt}\OperatorTok{[}\NormalTok{GG Msun}\OperatorTok{]}\NormalTok{;}
\FunctionTok{FullSimplify}\OperatorTok{[}\FunctionTok{T} \SpecialCharTok{{-}}\NormalTok{ T0 }\SpecialCharTok{/} \FunctionTok{Sqrt}\OperatorTok{[}\DecValTok{1} \SpecialCharTok{+}\NormalTok{ mMerc}\SpecialCharTok{/}\NormalTok{Msun}\OperatorTok{]]}
\CommentTok{(* Result: 0  [OK] *)}
\end{Highlighting}
\end{Shaded}

\textbf{Algebra {[}OK{]}.} Period decreases slightly because the
effective gravitational constant in the one-body problem is
\texttt{G(M+m)} not \texttt{GM}. This is a real kinematic effect on the
\textbf{sidereal period}.

\hypertarget{eqs.-1820-angular-velocity-and-perihelion-advance}{%
\subsubsection{§2.5 Eqs. (18)--(20) --- angular velocity and
``perihelion
advance''}\label{eqs.-1820-angular-velocity-and-perihelion-advance}}

\[
\omega = \omega_0\,(1 + m/M)^{1/2}             \qquad (18)
\] \[
\varphi = \omega T_0 = 2\pi\,(1 + m/M)^{1/2} \simeq 2\pi\,(1 + m/(2M))  \qquad (19)
\] \[
\boxed{\;\Delta\varphi = \varphi - \varphi_0 \simeq \pi m / M\;}  \qquad (20)
\]

\textbf{Wolfram MCP Taylor check:}

\begin{Shaded}
\begin{Highlighting}[]
\FunctionTok{Series}\OperatorTok{[}\DecValTok{2} \FunctionTok{Pi} \FunctionTok{Sqrt}\OperatorTok{[}\DecValTok{1} \SpecialCharTok{+} \FunctionTok{x}\OperatorTok{]} \SpecialCharTok{{-}} \DecValTok{2} \FunctionTok{Pi}\OperatorTok{,} \OperatorTok{\{}\FunctionTok{x}\OperatorTok{,} \DecValTok{0}\OperatorTok{,} \DecValTok{2}\OperatorTok{\}]}
\CommentTok{(* Result: Pi x {-} Pi x\^{}2/4 + O[x]\^{}3  [OK] leading term is π m/M *)}
\end{Highlighting}
\end{Shaded}

\textbf{Algebra {[}OK{]}.} Numerically for Mercury
(\texttt{m/M\ =\ 1.66\ ×\ 10⁻⁷}):
\texttt{Δφ\ =\ 5.21\ ×\ 10⁻⁷\ rad/orbit\ ×\ 415.2\ orbits/century\ ×\ (180·3600/π)}
\texttt{=\ 44.66\textquotesingle{}\textquotesingle{}/century}.
Wolfram-verified.

\textbf{{[}X{]} Physical-intuition status {[}for Tepper:{]}} here is the
load-bearing critique. \textbf{What Eq. (20) computes is the angle
Mercury sweeps in the unperturbed period \texttt{T\_0} at the perturbed
angular velocity \texttt{ω}.} That is not a perihelion advance. The
perihelion advance is the angular rotation of the orbit's major axis
between successive perihelia --- a property of the shape of the orbit,
not of the period or angular velocity.

For an inverse-square force the orbit is a closed ellipse (Bertrand's
theorem), so the angle between successive perihelia is exactly
\texttt{2π}, \textbf{regardless of \texttt{T} and \texttt{ω}}. Changing
\texttt{G} from \texttt{GM} to \texttt{G(M+m)} rescales the orbital
speed and period but leaves the orbit a closed ellipse. The
\texttt{πm/M} quantity Corda computes is a real kinematic effect, but it
is not the perihelion-advance quantity. The §3 derivation makes this
completely explicit --- see §2.7 below.

\hypertarget{eqs.-2131-cordas-own-apsidal-angle-formalism-3-of-paper}{%
\subsubsection{§2.6 Eqs. (21)--(31) --- Corda's own apsidal-angle
formalism (§3 of
paper)}\label{eqs.-2131-cordas-own-apsidal-angle-formalism-3-of-paper}}

Following Price \& Rush (1979), Corda writes the radial equation,
imposes angular-momentum conservation, Taylor-expands around
\texttt{r\_0}, and arrives at the \textbf{classical Bertrand-style
apsidal angle}:

\[
F_{c0}(r) = m\left(\ddot r - \omega_0^2 r\right)                  \qquad (21)
\] \[
J_0 = m r^2 \omega_0                                              \qquad (22)
\] \[
F_{c0}(r) = m\left(\ddot r - \frac{J_0^2}{m^2 r^3}\right)         \qquad (23)
\] \[
F_{c0}(r_0) = -\frac{J_0^2}{m r_0^3}                              \qquad (24)
\] After perturbing \texttt{x\ =\ r\ −\ r\_0} and series-expanding: \[
\ddot x + m^{-1}\!\left[-\frac{3 F_{c0}(r_0)}{r_0} - F'_{c0}(r_0)\right] x = 0  \qquad (27)
\] \[
T_0 = 2\pi\!\left(\frac{m}{-\,3 F_{c0}(r_0)/r_0 - F'_{c0}(r_0)}\right)^{1/2}  \qquad (28)
\] \[
\frac{\varphi_0}{2} = \pi\!\left(\frac{m}{-3F_{c0}/r_0 - F'_{c0}}\right)^{1/2}\!\frac{J_0}{m r_0^2}  \qquad (29)
\] Using Eqs. (24) and (29): \[
\boxed{\;\varphi_0 = 2\pi\left[\,3 + \frac{F'_{c0}(r_0)\,r_0}{F_{c0}(r_0)}\,\right]^{-1/2}\;}  \qquad (31)
\]

\textbf{Wolfram MCP for \texttt{F\ =\ −k/r²}:}

\begin{Shaded}
\begin{Highlighting}[]
\FunctionTok{ClearAll}\OperatorTok{[}\NormalTok{Fc}\OperatorTok{,} \FunctionTok{r}\OperatorTok{,}\NormalTok{ r0}\OperatorTok{,}\NormalTok{ kk}\OperatorTok{]}\NormalTok{;}
\NormalTok{Fc}\OperatorTok{[}\AttributeTok{r\_}\OperatorTok{]} \ExtensionTok{:=} \SpecialCharTok{{-}}\NormalTok{kk}\SpecialCharTok{/}\FunctionTok{r}\SpecialCharTok{\^{}}\DecValTok{2}\NormalTok{;}
\NormalTok{ratio }\ExtensionTok{=}\NormalTok{ (Fc\textquotesingle{}}\OperatorTok{[}\FunctionTok{r}\OperatorTok{]} \FunctionTok{r}\SpecialCharTok{/}\NormalTok{Fc}\OperatorTok{[}\FunctionTok{r}\OperatorTok{]}\NormalTok{) }\OtherTok{/.} \FunctionTok{r} \OtherTok{{-}\textgreater{}}\NormalTok{ r0;}
\NormalTok{phi0  }\ExtensionTok{=} \DecValTok{2} \FunctionTok{Pi} \SpecialCharTok{/} \FunctionTok{Sqrt}\OperatorTok{[}\DecValTok{3} \SpecialCharTok{+}\NormalTok{ ratio}\OperatorTok{]}\NormalTok{;}
\OperatorTok{\{}\FunctionTok{Simplify}\OperatorTok{[}\NormalTok{ratio}\OperatorTok{],} \FunctionTok{Simplify}\OperatorTok{[}\NormalTok{phi0}\OperatorTok{]\}}
\CommentTok{(* Result: \{{-}2, 2 Pi\}  [OK] For inverse{-}square: F\textquotesingle{}r/F = −2, φ\_0 = 2π, NO PRECESSION *)}
\end{Highlighting}
\end{Shaded}

\textbf{Algebra {[}OK{]} --- and this contradicts §2.} Corda derives the
apsidal-angle formula from scratch (Eq. 31), specialises to
inverse-square in his own text (line right after Eq. 31: \emph{``by
setting F\_\{c0\} = F\_G in Eq. (31)\ldots{} one finds φ\_0 = 2π, which
is exactly Eq. (7)''}), and explicitly notes that this is the
closed-orbit result. \textbf{For an inverse-square force, the apsidal
angle is 2π, the orbit is closed, and there is no precession.} Verified
by Wolfram MCP
(\href{https://en.wikipedia.org/wiki/Bertrand\%27s_theorem}{Bertrand's
theorem} at work).

\hypertarget{eqs.-3241-cordas-promotion-to-the-1mm-force}{%
\subsubsection{§2.7 Eqs. (32)--(41) --- Corda's ``promotion'' to the
(1+m/M) force}\label{eqs.-3241-cordas-promotion-to-the-1mm-force}}

To extend Eq. (31)'s apsidal-angle calculation to the reduced-mass case,
Corda makes the substitution: \[
F_{c0}(r) \;\to\; F_c(r) = \left(1 + \frac{m}{M}\right) F_{c0}(r)  \qquad (32)
\] \[
\omega_0 \to \omega                                                \qquad (33)
\] \[
J_0 \to J = m r^2 \omega                                           \qquad (34)
\] This is meant to mirror the §2 substitution
\texttt{M\ \textbackslash{}to\ M+m} inside the central force law. He
then writes: \[
F'_c(r_0) = \left(1 + \frac{m}{M}\right) F'_{c0}(r_0)              \qquad (36)
\] \[
\boxed{\;\frac{F'_c(r_0)}{F_c(r_0)} \;=\; \frac{F'_{c0}(r_0)}{F_{c0}(r_0)}\;}  \qquad (37)
\]

\textbf{Wolfram MCP check:}

\begin{Shaded}
\begin{Highlighting}[]
\FunctionTok{ClearAll}\OperatorTok{[}\NormalTok{Fc}\OperatorTok{,}\NormalTok{ Fc0}\OperatorTok{,} \FunctionTok{r}\OperatorTok{,}\NormalTok{ r0}\OperatorTok{,}\NormalTok{ mMerc}\OperatorTok{,}\NormalTok{ Msun}\OperatorTok{]}\NormalTok{;}
\NormalTok{Fc}\OperatorTok{[}\AttributeTok{r\_}\OperatorTok{]} \ExtensionTok{:=}\NormalTok{ (}\DecValTok{1} \SpecialCharTok{+}\NormalTok{ mMerc}\SpecialCharTok{/}\NormalTok{Msun) Fc0}\OperatorTok{[}\FunctionTok{r}\OperatorTok{]}\NormalTok{;}
\FunctionTok{FullSimplify}\OperatorTok{[}\NormalTok{Fc\textquotesingle{}}\OperatorTok{[}\NormalTok{r0}\OperatorTok{]}\SpecialCharTok{/}\NormalTok{Fc}\OperatorTok{[}\NormalTok{r0}\OperatorTok{]} \SpecialCharTok{{-}}\NormalTok{ Fc0\textquotesingle{}}\OperatorTok{[}\NormalTok{r0}\OperatorTok{]}\SpecialCharTok{/}\NormalTok{Fc0}\OperatorTok{[}\NormalTok{r0}\OperatorTok{]]}
\CommentTok{(* Result: 0  [OK] — the (1+m/M) factor cancels in F\textquotesingle{}/F *)}
\end{Highlighting}
\end{Shaded}

\textbf{Algebra {[}OK{]}, and {[}X{]} here is where Corda's apparatus
contradicts his own §2 conclusion.} The apsidal-angle formula Eq. (31)
depends \emph{only} on the ratio \texttt{F\textquotesingle{}/F}. The Eq.
(37) Wolfram-verified identity says that constant rescaling of
\texttt{F\_\{c0\}} leaves \texttt{F\textquotesingle{}/F} unchanged.
Therefore the apsidal angle for the rescaled force
\texttt{F\_c\ =\ (1+m/M)\ F\_\{c0\}} is \emph{the same} as for the
unscaled force \texttt{F\_\{c0\}} --- which is \texttt{2π} (Eq. 31 with
\texttt{F\_\{c0\}\ =\ F\_G\ =\ −k/r²}). \textbf{The ``perturbed'' force
still gives a closed orbit. No precession.}

Corda's subsequent Eqs. (38)--(41) derive
\texttt{T\ =\ T\_0(1+m/M)\^{}(−1/2)} again from the period formula (28),
then \texttt{ω\ =\ ω\_0(1+m/M)\^{}(1/2)}, then
\texttt{φ\ =\ ω\ T\_0\ ≈\ 2π(1\ +\ m/(2M))} --- the same \texttt{πm/M}
as §2. \textbf{But this is exactly the §2 mistake repeated: a period
change rephrased as a ``perihelion advance.''} Corda's own Eq. (31),
applied to his own Eq. (32) force, \emph{contradicts} the §3 conclusion.

\hypertarget{eqs.-4257-reduced-mass-kepler-derivation}{%
\subsubsection{§2.8 Eqs. (42)--(57) --- reduced-mass Kepler
derivation}\label{eqs.-4257-reduced-mass-kepler-derivation}}

Corda's §4 (``Third Kepler's law'') works out the standard
centre-of-mass + reduced-mass two-body decomposition: \[
\vec R = \frac{m\vec r_m + M\vec r_M}{M + m},\qquad
\vec r = \vec r_m - \vec r_M                                       \qquad (43)
\] \[
J = \mu r^2 \omega = 2\mu\,\dot A                                  \qquad (45)
\] \[
T = 2\pi\!\sqrt{\frac{a^3}{G M_T}},\quad M_T = M + m              \qquad (53)
\] \[
\frac{a^3}{T^2} = \frac{G(M+m)}{4\pi^2}                            \qquad (54)
\] \[
T = T_0 (1 + m/M)^{-1/2}                                           \qquad (57)
\]

\textbf{Algebra {[}OK{]}} --- standard textbook treatment (Goldstein
§3.7, MTW §25). Lands at the same \texttt{T\ =\ T\_0(1+m/M)\^{}(−1/2)}
as Eqs. (17) and (39).

\textbf{Physical-intuition status {[}for Tepper:{]}} §4's derivation is
impeccable as a derivation of the \emph{period correction} due to the
two-body kinematics. That's what the reduced-mass treatment gives. But
Corda then writes Eq. (59) \texttt{φ\ =\ ωT\_0\ ≈\ 2π(1\ +\ m/(2M))} and
calls the excess a ``perihelion advance'' --- same conflation as §2.7.

\hypertarget{numerical-summary-of-cordas-three-approaches-234-all-give-the-same-ux3b4ux3c6}{%
\subsubsection{§2.9 Numerical summary of Corda's three approaches
(§2/§3/§4 all give the same
Δφ)}\label{numerical-summary-of-cordas-three-approaches-234-all-give-the-same-ux3b4ux3c6}}

\begin{Shaded}
\begin{Highlighting}[]
\FunctionTok{ClearAll}\OperatorTok{[}\NormalTok{GG}\OperatorTok{,}\NormalTok{ Msun}\OperatorTok{,}\NormalTok{ mMerc}\OperatorTok{,}\NormalTok{ aMerc}\OperatorTok{,}\NormalTok{ Torb}\OperatorTok{,}\NormalTok{ orbitsPerCentury}\OperatorTok{,}\NormalTok{ radToArcsec}\OperatorTok{,}\NormalTok{ dphiNewton}\OperatorTok{]}\NormalTok{;}
\NormalTok{GG }\ExtensionTok{=} \FloatTok{6.6743}\SpecialCharTok{*\^{}{-}}\DecValTok{11}\NormalTok{; Msun }\ExtensionTok{=} \FloatTok{1.98892}\SpecialCharTok{*\^{}}\DecValTok{30}\NormalTok{; mMerc }\ExtensionTok{=} \FloatTok{3.3011}\SpecialCharTok{*\^{}}\DecValTok{23}\NormalTok{;}
\NormalTok{aMerc }\ExtensionTok{=} \FloatTok{5.7909}\SpecialCharTok{*\^{}}\DecValTok{10}\NormalTok{; Torb }\ExtensionTok{=} \FloatTok{87.969}\SpecialCharTok{*}\FloatTok{86400.}\NormalTok{;}
\NormalTok{orbitsPerCentury }\ExtensionTok{=} \DecValTok{100} \SpecialCharTok{*} \FloatTok{365.25} \SpecialCharTok{*} \FloatTok{86400.} \SpecialCharTok{/}\NormalTok{ Torb;}
\NormalTok{radToArcsec }\ExtensionTok{=} \DecValTok{180}\SpecialCharTok{/}\FunctionTok{Pi} \SpecialCharTok{*} \FloatTok{3600.}\NormalTok{;}
\NormalTok{dphiNewton }\ExtensionTok{=} \FunctionTok{Pi}\NormalTok{ mMerc }\SpecialCharTok{/}\NormalTok{ Msun;}
\FunctionTok{Print}\OperatorTok{[}\StringTok{"Δφ = πm/M per orbit  = "}\OperatorTok{,} \FunctionTok{N}\OperatorTok{[}\NormalTok{dphiNewton}\OperatorTok{],} \StringTok{" rad"}\OperatorTok{]}\NormalTok{;}
\FunctionTok{Print}\OperatorTok{[}\StringTok{"Δφ = πm/M per century = "}\OperatorTok{,} \FunctionTok{N}\OperatorTok{[}\NormalTok{dphiNewton }\SpecialCharTok{*}\NormalTok{ orbitsPerCentury }\SpecialCharTok{*}\NormalTok{ radToArcsec}\OperatorTok{],} \StringTok{" arcsec"}\OperatorTok{]}\NormalTok{;}
\CommentTok{(* Result:                                                                       *)}
\CommentTok{(*   Δφ = πm/M per orbit  = 5.21424×10⁻⁷ rad                                     *)}
\CommentTok{(*   Δφ = πm/M per century = 44.6557 arcsec  ← Corda\textquotesingle{}s reported 44.39\textquotesingle{}\textquotesingle{}/c        *)}
\end{Highlighting}
\end{Shaded}

\textbf{The numerical match to observed
\texttt{43\textquotesingle{}\textquotesingle{}/century} is correct
arithmetic on a quantity that --- per §2.6 and §2.7 above --- is }not**
a perihelion advance.**

\hypertarget{breakdown-for-venus-and-earth-the-papers-self-refutation}{%
\subsubsection{§2.10 §5 --- Breakdown for Venus and Earth (the paper's
self-refutation)}\label{breakdown-for-venus-and-earth-the-papers-self-refutation}}

Corda himself shows the \texttt{πm/M} formula fails for non-Mercury
planets:

\begin{Shaded}
\begin{Highlighting}[]
\FunctionTok{ClearAll}\OperatorTok{[}\NormalTok{mVenus}\OperatorTok{,}\NormalTok{ mEarth}\OperatorTok{,}\NormalTok{ TVenus}\OperatorTok{,}\NormalTok{ TEarth}\OperatorTok{,}\NormalTok{ orbitsVenus}\OperatorTok{,}\NormalTok{ orbitsEarth}\OperatorTok{]}\NormalTok{;}
\NormalTok{mVenus }\ExtensionTok{=} \FloatTok{4.87}\SpecialCharTok{*\^{}}\DecValTok{24}\NormalTok{; TVenus }\ExtensionTok{=} \FloatTok{224.701}\SpecialCharTok{*}\FloatTok{86400.}\NormalTok{;}
\NormalTok{mEarth }\ExtensionTok{=} \FloatTok{5.97}\SpecialCharTok{*\^{}}\DecValTok{24}\NormalTok{; TEarth }\ExtensionTok{=} \FloatTok{365.256}\SpecialCharTok{*}\FloatTok{86400.}\NormalTok{;}
\NormalTok{orbitsVenus }\ExtensionTok{=} \DecValTok{100}\SpecialCharTok{*}\FloatTok{365.25}\SpecialCharTok{*}\FloatTok{86400.}\SpecialCharTok{/}\NormalTok{TVenus;}
\NormalTok{orbitsEarth }\ExtensionTok{=} \DecValTok{100}\SpecialCharTok{*}\FloatTok{365.25}\SpecialCharTok{*}\FloatTok{86400.}\SpecialCharTok{/}\NormalTok{TEarth;}
\OperatorTok{\{}\StringTok{"Venus"}\OperatorTok{,}  \FunctionTok{N}\OperatorTok{[}\FunctionTok{Pi}\NormalTok{ mVenus}\SpecialCharTok{/}\NormalTok{Msun }\SpecialCharTok{*}\NormalTok{ orbitsVenus }\SpecialCharTok{*}\NormalTok{ radToArcsec}\OperatorTok{],} \StringTok{"vs observed 8.62\textquotesingle{}\textquotesingle{}/c"}\OperatorTok{\}}
\OperatorTok{\{}\StringTok{"Earth"}\OperatorTok{,}  \FunctionTok{N}\OperatorTok{[}\FunctionTok{Pi}\NormalTok{ mEarth}\SpecialCharTok{/}\NormalTok{Msun }\SpecialCharTok{*}\NormalTok{ orbitsEarth }\SpecialCharTok{*}\NormalTok{ radToArcsec}\OperatorTok{],} \StringTok{"vs observed 3.83\textquotesingle{}\textquotesingle{}/c"}\OperatorTok{\}}
\CommentTok{(* Results:                                                                      *)}
\CommentTok{(*   Venus  formula prediction: 257.91 \textquotesingle{}\textquotesingle{}/c   vs observed  8.62\textquotesingle{}\textquotesingle{}/c  →  30× too big *)}
\CommentTok{(*   Earth  formula prediction: 194.50 \textquotesingle{}\textquotesingle{}/c   vs observed  3.83\textquotesingle{}\textquotesingle{}/c  →  51× too big *)}
\end{Highlighting}
\end{Shaded}

\begin{longtable}[]{@{}llll@{}}
\toprule\noalign{}
Planet & \texttt{πm/M} (Corda's formula) & Observed residual &
Discrepancy \\
\midrule\noalign{}
\endhead
\bottomrule\noalign{}
\endlastfoot
Mercury & \texttt{44.66\textquotesingle{}\textquotesingle{}/c} &
\texttt{≈\ 43\textquotesingle{}\textquotesingle{}/c} & {[}OK{]} 1.04× ←
``agreement'' \\
Venus & \texttt{257.91\textquotesingle{}\textquotesingle{}/c} &
\texttt{≈\ 8.62\textquotesingle{}\textquotesingle{}/c} & {[}X{]}
\textbf{30× too big} \\
Earth & \texttt{194.50\textquotesingle{}\textquotesingle{}/c} &
\texttt{≈\ 3.83\textquotesingle{}\textquotesingle{}/c} & {[}X{]}
\textbf{51× too big} \\
\end{longtable}

\textbf{A formula that's right by 4\% for one planet and wrong by
factors of 30--50 for others is not a physics result --- it is a
numerical coincidence at one particular \texttt{m/M} value.} Corda
himself acknowledges the breakdown (§5) but rationalises it with n-body
effects (``Venus has Mercury as an interior planet''), which does not
explain why a strictly two-body formula would be valid for one two-body
system and invalid for others. \textbf{The honest reading: \texttt{πm/M}
is not the perihelion advance for any planet, including Mercury.}

\hypertarget{eqs.-6067-gravitational-time-dilation}{%
\subsubsection{§2.11 Eqs. (60)--(67) --- gravitational time
dilation}\label{eqs.-6067-gravitational-time-dilation}}

Corda's §6 applies Schwarzschild weak-field corrections. Standard
relations: \[
t_g = \sqrt{g_{00}(r_0)}\, t_l                                     \qquad (60)
\] Isotropic-coordinate weak-field Schwarzschild line element: \[
ds^2 = \left(1 - \frac{2GM}{r c^2}\right)(c\,dt)^2
       - \left(1 + \frac{2GM}{r c^2}\right)(dr^2 + r^2 d\theta^2 + r^2 \sin^2\theta\,d\phi^2)
                                                                   \qquad (61)
\] \[
t_g = \sqrt{1 - r_g/r_0}\, t_l \;\approx\; \left(1 - \tfrac{1}{2}\tfrac{r_g}{r_0}\right) t_l   \qquad (62)
\] where \texttt{r\_g\ =\ 2GM/c²} is the Sun's gravitational radius
(\texttt{≈\ 2.95\ km}). Corda also asserts a corresponding
radial-distance contraction: \texttt{r₀\ →\ r₀(1\ −\ ½\ r\_g/r₀)} (Eq.
63), arguing from \texttt{r\ =\ ct}.

\textbf{{[}partial{]} Physical-intuition status {[}for Tepper:{]}} the
time-dilation factor (Eq. 62) is standard GR weak-field. The
radial-distance contraction in Eq. (63) is unusual --- in isotropic
coordinates the radial coordinate \texttt{r} of a point at Schwarzschild
radius \texttt{R} is \texttt{r\ ≈\ R(1\ −\ GM/(c²R))} (a small
isotropic-to-areal coordinate shift), so there \emph{is} a
\texttt{½\ r\_g/r₀} correction when comparing isotropic to areal radial
coordinates, but treating this as ``the CNO sees a different distance
because \texttt{r\ =\ ct}'' is non-standard. Worth flagging but the
final-result coefficient in Eq. (67) is unaffected to first order.

Substituting Eq. (63) into Eq. (17): \[
T_F = \frac{2\pi r_0^{3/2}\left(1 - \tfrac{1}{2}\tfrac{r_g}{r_0}\right)^{3/2}(1 + m/M)^{-1/2}}{\sqrt{GM}}  \qquad (64)
\] Computing \texttt{ω\_F\ =\ 2π/T\_F}: \[
\omega_F \simeq \omega\left(1 - \tfrac{1}{2}\tfrac{r_g}{r_0}\right)^{-3/2}\!\left(1 + \tfrac{1}{2}\tfrac{r_g}{r_0}\right)   \qquad (65)
\] \[
\varphi_F \simeq 2\pi\left(1 + \tfrac{5}{4}\tfrac{r_g}{r_0}\right)   \qquad (66)
\] \[
\boxed{\;\Delta\varphi_F \;\simeq\; \tfrac{5\pi}{2}\,\tfrac{r_g}{r_0}\;}  \qquad (67)
\]

\textbf{Wolfram MCP Taylor check:}

\begin{Shaded}
\begin{Highlighting}[]
\FunctionTok{Series}\OperatorTok{[}\DecValTok{2} \FunctionTok{Pi}\NormalTok{ (}\DecValTok{1} \SpecialCharTok{{-}} \FunctionTok{x}\SpecialCharTok{/}\DecValTok{2}\NormalTok{)}\SpecialCharTok{\^{}}\NormalTok{(}\SpecialCharTok{{-}}\DecValTok{3}\SpecialCharTok{/}\DecValTok{2}\NormalTok{) (}\DecValTok{1} \SpecialCharTok{+} \FunctionTok{x}\SpecialCharTok{/}\DecValTok{2}\NormalTok{)}\OperatorTok{,} \OperatorTok{\{}\FunctionTok{x}\OperatorTok{,} \DecValTok{0}\OperatorTok{,} \DecValTok{1}\OperatorTok{\}]}
\CommentTok{(* Result: 2 Pi + (5 Pi/2) x + O[x]\^{}2   [OK] leading is (5/2)π r\_g/r\_0 *)}
\end{Highlighting}
\end{Shaded}

\textbf{Algebra {[}OK{]}} given the substitution in Eq. (63). This is
roughly two-thirds of the GR answer.

\hypertarget{eqs.-6872-rotational-time-dilation-langevin}{%
\subsubsection{§2.12 Eqs. (68)--(72) --- rotational time dilation
(Langevin)}\label{eqs.-6872-rotational-time-dilation-langevin}}

Standard SR derivation of the Langevin line element for a rotating
frame: \[
ds^2 = c^2 dt^2 - dr^2 - r^2 d\phi^2 - dz^2                       \qquad (68)
\] \[
\phi = \phi' + \omega t'                                           \qquad (69)
\] \[
ds^2 = \left(1 - \tfrac{r^2\omega^2}{c^2}\right)c^2 dt^2 - 2\omega r^2 d\phi' dt' - dr^2 - r^2 d\phi'^2 - dz'^2   \qquad (70)
\] \[
d\tau = \sqrt{1 - r^2\omega^2/c^2}\,dt \;\simeq\; \left(1 - \tfrac{1}{2}\tfrac{r^2\omega^2}{c^2}\right) dt   \qquad (71)
\] At \texttt{r\ =\ r\_0}, \texttt{ω\ =\ ω\_0}: \[
\tau \simeq \left(1 - \tfrac{1}{2}\tfrac{r_0^2\omega_0^2}{c^2}\right) t   \qquad (72)
\]

\textbf{Algebra {[}OK{]}} --- textbook special relativity.
\textbf{Physical-intuition status:} Corda invokes Einstein's Equivalence
Principle to argue that the centrifugal field in the rotating frame can
be treated ``as if'' it were a gravitational field. This is the
conceptual move that permits combining gravitational + rotational TD
additively. \textbf{{[}for Tepper:{]}} the EEP application here is
standard; the question is whether the Langevin correction, applied to a
Keplerian orbital period, is the right way to recover the GR perihelion
advance. The §8 result suggests yes (in the circular limit), and this
matches the standard PPN result for the gravitational-redshift + Doppler
combination.

\hypertarget{eqs.-7377-combined-correction}{%
\subsubsection{§2.13 Eqs. (73)--(77) --- combined
correction}\label{eqs.-7377-combined-correction}}

Using Eqs. (5)--(6) to express the rotational TD in terms of
\texttt{r\_g}: \[
\tau = \sqrt{1 - \tfrac{1}{2}\tfrac{r_g}{r_0}}\, t
     \simeq \left(1 - \tfrac{1}{4}\tfrac{r_g}{r_0}\right) t        \qquad (73)
\] The CNO replaces
\texttt{T\_F\ \textbackslash{}to\ T\_T\ =\ T\_F\ (1\ -\ \textbackslash{}tfrac\{1\}\{2\}\ r\_g/r\_0)\^{}\{1/2\}}
and: \[
T_T = T_0 (1 + m/M)^{-1/2}\,(1 - \tfrac{1}{2} r_g/r_0)^2           \qquad (74)
\] \[
\omega_F = 2\pi/T_T \simeq \omega\,(1 - \tfrac{1}{2} r_g/r_0)^{-2} \simeq \omega\,(1 + \tfrac{3}{2}\, r_g/r_0)   \qquad (75)
\] \[
\varphi_F = \omega_F T_0 (1 - m/(2M)) \simeq 2\pi\left(1 + \tfrac{3}{2}\,\tfrac{r_g}{r_0}\right)  \qquad (76)
\] \[
\boxed{\;\Delta\varphi_F \;\simeq\; 3\pi\,\tfrac{r_g}{r_0}\;}  \qquad (77)
\]

\begin{quote}
\textbf{OCR/marker-pdf note.} The conversion of Eq. (75) shows the last
rotational factor as \texttt{(1\ -\ r\_g/(2r\_0))\^{}\{+1/2\}}, which
would give a leading coefficient of \texttt{1} (not \texttt{3/2}).
Reading instead as \texttt{(1\ -\ r\_g/(2r\_0))\^{}\{-1/2\}} (consistent
with \texttt{ω\ =\ 2π/T} and the
\texttt{T\_T\ =\ T\_F\ (1\ -\ r\_g/(2r\_0))\^{}\{1/2\}} substitution)
recovers the \texttt{3/2}. Either Eq. (75) in the paper has a typo or
the converted markdown does. The final Eq. (77) is correct and matches
the standard PPN gravitational-redshift + Doppler combination.
\end{quote}

\textbf{Wolfram MCP Taylor check (assuming \texttt{−1/2}):}

\begin{Shaded}
\begin{Highlighting}[]
\FunctionTok{Series}\OperatorTok{[}\NormalTok{(}\DecValTok{1} \SpecialCharTok{{-}} \FunctionTok{x}\SpecialCharTok{/}\DecValTok{2}\NormalTok{)}\SpecialCharTok{\^{}}\NormalTok{(}\SpecialCharTok{{-}}\DecValTok{3}\SpecialCharTok{/}\DecValTok{2}\NormalTok{) (}\DecValTok{1} \SpecialCharTok{+} \FunctionTok{x}\SpecialCharTok{/}\DecValTok{2}\NormalTok{) (}\DecValTok{1} \SpecialCharTok{{-}} \FunctionTok{x}\SpecialCharTok{/}\DecValTok{2}\NormalTok{)}\SpecialCharTok{\^{}}\NormalTok{(}\SpecialCharTok{{-}}\DecValTok{1}\SpecialCharTok{/}\DecValTok{2}\NormalTok{)}\OperatorTok{,} \OperatorTok{\{}\FunctionTok{x}\OperatorTok{,} \DecValTok{0}\OperatorTok{,} \DecValTok{1}\OperatorTok{\}]}
\CommentTok{(* Result: 1 + 3x/2 + O[x]\^{}2   [OK] leading is (3/2) r\_g/r\_0 *)}
\end{Highlighting}
\end{Shaded}

\textbf{Algebra {[}OK{]}} (modulo the suspected typo in Eq. 75).

\hypertarget{comparison-of-cordas-eq.-77-to-standard-gr}{%
\subsubsection{§2.14 Comparison of Corda's Eq. (77) to standard
GR}\label{comparison-of-cordas-eq.-77-to-standard-gr}}

Standard GR (Schwarzschild geodesic, MTW §40.5): \[
\Delta\varphi_{\rm GR} = \frac{6\pi GM}{c^2 a (1-e^2)} = \frac{3\pi r_g}{a(1-e^2)}
\] For Mercury (\texttt{e\ =\ 0.20563},
\texttt{1\ -\ e\^{}2\ ≈\ 0.9577}):
\texttt{Δφ\_GR\ =\ 42.99\textquotesingle{}\textquotesingle{}/century}.

For the \textbf{circular} limit (\texttt{e\ =\ 0}, \texttt{a\ =\ r\_0}):
\texttt{Δφ\_GR\ =\ 3π\ r\_g/r\_0}. \textbf{Corda's Eq. (77) reproduces
the circular-orbit limit of GR.}

\begin{Shaded}
\begin{Highlighting}[]
\FunctionTok{ClearAll}\OperatorTok{[}\NormalTok{GG}\OperatorTok{,}\NormalTok{ Msun}\OperatorTok{,}\NormalTok{ aMerc}\OperatorTok{,}\NormalTok{ eMerc}\OperatorTok{,}\NormalTok{ cc}\OperatorTok{,}\NormalTok{ rg}\OperatorTok{,}\NormalTok{ orbitsPerCentury}\OperatorTok{,}\NormalTok{ radToArcsec}\OperatorTok{]}\NormalTok{;}
\NormalTok{GG }\ExtensionTok{=} \FloatTok{6.6743}\SpecialCharTok{*\^{}{-}}\DecValTok{11}\NormalTok{; Msun }\ExtensionTok{=} \FloatTok{1.98892}\SpecialCharTok{*\^{}}\DecValTok{30}\NormalTok{; aMerc }\ExtensionTok{=} \FloatTok{5.7909}\SpecialCharTok{*\^{}}\DecValTok{10}\NormalTok{;}
\NormalTok{eMerc }\ExtensionTok{=} \FloatTok{0.20563}\NormalTok{; cc }\ExtensionTok{=} \FloatTok{299792458.}\NormalTok{; rg }\ExtensionTok{=} \DecValTok{2}\NormalTok{ GG Msun }\SpecialCharTok{/}\NormalTok{ cc}\SpecialCharTok{\^{}}\DecValTok{2}\NormalTok{;}
\NormalTok{orbitsPerCentury }\ExtensionTok{=} \DecValTok{100} \SpecialCharTok{*} \FloatTok{365.25} \SpecialCharTok{*} \FloatTok{86400.} \SpecialCharTok{/}\NormalTok{ (}\FloatTok{87.969} \SpecialCharTok{*} \FloatTok{86400.}\NormalTok{);}
\NormalTok{radToArcsec }\ExtensionTok{=} \DecValTok{180}\SpecialCharTok{/}\FunctionTok{Pi} \SpecialCharTok{*} \FloatTok{3600.}\NormalTok{;}
\OperatorTok{\{}\StringTok{"Corda §8 (3π r\_g/a, circular)"}\OperatorTok{,}  \DecValTok{3} \FunctionTok{Pi}\NormalTok{ rg}\SpecialCharTok{/}\NormalTok{aMerc }\SpecialCharTok{*}\NormalTok{ orbitsPerCentury }\SpecialCharTok{*}\NormalTok{ radToArcsec}\OperatorTok{,}
 \StringTok{"GR Mercury  (3π r\_g/[a(1{-}e²)])"}\OperatorTok{,} \DecValTok{3} \FunctionTok{Pi}\NormalTok{ rg}\SpecialCharTok{/}\NormalTok{(aMerc (}\DecValTok{1} \SpecialCharTok{{-}}\NormalTok{ eMerc}\SpecialCharTok{\^{}}\DecValTok{2}\NormalTok{)) }\SpecialCharTok{*}\NormalTok{ orbitsPerCentury }\SpecialCharTok{*}\NormalTok{ radToArcsec}\OperatorTok{\}}
\CommentTok{(* Results:                                                                  *)}
\CommentTok{(*   Corda §8 circular:  41.17\textquotesingle{}\textquotesingle{}/c                                            *)}
\CommentTok{(*   GR Mercury (e=0.21): 42.99\textquotesingle{}\textquotesingle{}/c                                           *)}
\end{Highlighting}
\end{Shaded}

Corda's §§6--8 derivation is missing the \texttt{1/(1-e²)} eccentricity
factor. For Mercury this is a \texttt{4.4\%} shortfall
(\texttt{−1.82\textquotesingle{}\textquotesingle{}/century}).
\textbf{{[}for Tepper:{]}} this discrepancy is significant for Mercury
specifically because of its high eccentricity; for Venus and Earth (more
circular) the discrepancy shrinks. A true derivation of the GR result
needs the eccentricity from the genuine orbit (e.g.\textasciitilde via
the apsidal-angle formula or the relativistic correction to the orbit
equation).

\begin{center}\rule{0.5\linewidth}{0.5pt}\end{center}

\hypertarget{the-load-bearing-critique-three-findings}{%
\subsection{§3. The load-bearing critique --- three
findings}\label{the-load-bearing-critique-three-findings}}

\hypertarget{finding-1-algebra-driven.-bertrands-theorem-rules-out-24}{%
\subsubsection{Finding 1 (algebra-driven). Bertrand's theorem rules out
§2--§4}\label{finding-1-algebra-driven.-bertrands-theorem-rules-out-24}}

Corda's apsidal-angle formula (Eq. 31), applied to \emph{any}
inverse-square force --- including the \texttt{(1+m/M)}-rescaled version
of Eq. (32) --- gives \texttt{φ\_0\ =\ 2π} exactly: -
\texttt{F\ =\ −k/r²} ⇒
\texttt{F\textquotesingle{}(r)\ r\ /\ F(r)\ =\ −2} ⇒
\texttt{φ\_0\ =\ 2π/√(3\ −\ 2)\ =\ 2π}. {[}OK{]} -
\texttt{F\ =\ (1+m/M)(−k/r²)} ⇒ same
\texttt{F\textquotesingle{}/F\ =\ −2/r} ⇒ same \texttt{φ\_0\ =\ 2π}.
{[}OK{]}

This is Bertrand's theorem: of all central potentials, \emph{only}
\texttt{1/r} (Kepler / Newton) and \texttt{r²} (harmonic) give closed
bound orbits universally. Inverse-square forces (Newton's gravity, Eq.
9) produce closed ellipses with no precession at any energy or angular
momentum. Constant rescaling of an inverse-square force does not change
this --- it just changes the orbital period.

\textbf{Therefore Eq. (20)'s \texttt{Δφ\ =\ πm/M} is not an apsidal
precession.} It is the angle by which Mercury's \emph{position} differs
from \texttt{2π} when measured at time \texttt{T\_0} (the unperturbed
period) under the perturbed angular velocity \texttt{ω}. After the true
perturbed period \texttt{T\ =\ T\_0(1+m/M)\^{}(−1/2)} has elapsed,
Mercury has swept exactly \texttt{2π} again and returned to perihelion.
\textbf{No precession.}

\hypertarget{finding-2-empirics-driven.-the-5-self-refutation}{%
\subsubsection{Finding 2 (empirics-driven). The §5
self-refutation}\label{finding-2-empirics-driven.-the-5-self-refutation}}

Applying Corda's own formula \texttt{Δφ\ =\ πm/M\ ×\ (orbits/century)}
to Venus and Earth:

\begin{longtable}[]{@{}
  >{\raggedright\arraybackslash}p{(\columnwidth - 6\tabcolsep) * \real{0.2500}}
  >{\raggedright\arraybackslash}p{(\columnwidth - 6\tabcolsep) * \real{0.2500}}
  >{\raggedright\arraybackslash}p{(\columnwidth - 6\tabcolsep) * \real{0.2500}}
  >{\raggedright\arraybackslash}p{(\columnwidth - 6\tabcolsep) * \real{0.2500}}@{}}
\toprule\noalign{}
\begin{minipage}[b]{\linewidth}\raggedright
Planet
\end{minipage} & \begin{minipage}[b]{\linewidth}\raggedright
\texttt{m\ ×\ orbits/c\ ×\ radToArcsec\ /\ Msun\ ×\ π}
\end{minipage} & \begin{minipage}[b]{\linewidth}\raggedright
Observed
\end{minipage} & \begin{minipage}[b]{\linewidth}\raggedright
Discrepancy
\end{minipage} \\
\midrule\noalign{}
\endhead
\bottomrule\noalign{}
\endlastfoot
Mercury & \texttt{44.66\textquotesingle{}\textquotesingle{}/c} &
\texttt{43\textquotesingle{}\textquotesingle{}/c} & {[}OK{]} \\
Venus & \texttt{257.91\textquotesingle{}\textquotesingle{}/c} &
\texttt{8.62\textquotesingle{}\textquotesingle{}/c} & {[}X{]}
\textbf{\texttt{30×}} \\
Earth & \texttt{194.50\textquotesingle{}\textquotesingle{}/c} &
\texttt{3.83\textquotesingle{}\textquotesingle{}/c} & {[}X{]}
\textbf{\texttt{51×}} \\
\end{longtable}

Corda's response (``n-body effects, Mercury has no interior planets'')
doesn't work as a defence: - The formula \texttt{Δφ\ =\ πm/M} makes no
reference to n-body effects. - The derivation in §§2/3/4 is strictly
two-body throughout (no interior planets enter). - If the formula were
the correct two-body result, it would apply to any two-body sub-problem
and the n-body corrections would be smaller perturbations --- not
corrections of \texttt{factor\ 30–50}.

\textbf{{[}for Tepper:{]}} the strongest available evidence that
\texttt{πm/M} is not the correct two-body perihelion advance is that the
formula is empirically wrong for every planet other than Mercury. The
Mercury ``agreement'' is coincidence at \texttt{m/M\ ≈\ 1.66\ ×\ 10⁻⁷}
happening to give a number near
\texttt{43\textquotesingle{}\textquotesingle{}/c} once multiplied by
Mercury's \texttt{415.2\ orbits/century}. Venus and Earth have different
\texttt{(m/M,\ orbits/c)} products and give wildly wrong predictions.

\hypertarget{finding-3-formal-derivation-driven.-68-reproduce-only-the-e0-gr-limit}{%
\subsubsection{Finding 3 (formal-derivation-driven). §§6--8 reproduce
only the e=0 GR
limit}\label{finding-3-formal-derivation-driven.-68-reproduce-only-the-e0-gr-limit}}

Corda's Eq. (77) \texttt{Δφ\_F\ =\ 3π\ r\_g/r\_0} is correct algebra
(modulo the §2.13 typo) but recovers only the circular-orbit limit of
the true GR result \texttt{3π\ r\_g/{[}a(1−e²){]}}. For Mercury (e =
0.21) this is a 4.4\% shortfall
(\texttt{−1.82\textquotesingle{}\textquotesingle{}/century}).

\textbf{{[}for Tepper:{]}} §§6--8 is a real first-order GR derivation
route; what it gets right is the ``PPN sum of gravitational + Doppler
time dilation'' result, which is well-known in GR pedagogy. What it
doesn't capture is the eccentricity dependence, which makes it
numerically equivalent to GR only for nearly-circular orbits. For
Mercury, the eccentricity correction matters; for Venus and Earth
(smaller \texttt{e}), \texttt{3π\ r\_g/a} is much closer to the full GR
answer.

\textbf{The two §§6--8 results work in their domain.} The trouble is
purely the combination with §§2--4. Corda writes (§9) that ``Newtonian +
corrections matches GR.'' The honest reading: §§6--8 \emph{alone}
matches GR (in the circular limit); §§2--4 contribute a fictitious
\texttt{πm/M} that doesn't correspond to a real precession; the two
together don't add up to GR --- the \texttt{πm/M} is just spurious.

\begin{center}\rule{0.5\linewidth}{0.5pt}\end{center}

\hypertarget{implication-for-teppers-relativistic-newtonian-theory}{%
\subsection{\texorpdfstring{§4. Implication for Tepper's
\emph{Relativistic Newtonian
Theory}}{§4. Implication for Tepper's Relativistic Newtonian Theory}}\label{implication-for-teppers-relativistic-newtonian-theory}}

Gill's \emph{Relativistic Newtonian Theory} extends Corda's framework by
adding the dual-theory \texttt{(1\ +\ V/(mc²))} factor to the force law:

\[
a_m = -\frac{GM}{r^2}\left(1 - \frac{GM}{c^2 r}\right)\hat u_r   \qquad (\text{Eq. h4})
\]

(transcribed from PR
\href{https://github.com/temoTxt/PyPhysics/pull/83}{\#83}'s
\href{../Equation_Verification/Dual_Newtonian_Theory.md}{\texttt{Equation\_Verification/Dual\_Newtonian\_Theory.md}};
all algebra Wolfram-MCP verified).

\hypertarget{what-pr-83-computed-corda-style-template}{%
\subsubsection{§4.1 What PR \#83 computed (Corda-style
template)}\label{what-pr-83-computed-corda-style-template}}

PR \#83
\href{../Mercury_Perihelion/03_numerical_predictions.md}{\texttt{03\_numerical\_predictions.md}}
followed Corda's heuristic:

\[
\Delta\varphi_d = 2\pi\!\left[\left(1 + \tfrac{m}{M}\right)
                  \left(1 - \tfrac{G(M^2 + m^2)}{(M+m) c^2 r_0}\right)\right]^{1/2} - 2\pi
\]

with first-order Taylor expansion:

\[
\Delta\varphi_{d_1} \;\simeq\; 2\pi\!\left[\frac{m}{2M} - \frac{GM(1+m/M)}{2 c^2 r_0}\right]
\]

This decomposes as:

\begin{longtable}[]{@{}
  >{\raggedright\arraybackslash}p{(\columnwidth - 4\tabcolsep) * \real{0.3333}}
  >{\raggedright\arraybackslash}p{(\columnwidth - 4\tabcolsep) * \real{0.3333}}
  >{\raggedright\arraybackslash}p{(\columnwidth - 4\tabcolsep) * \real{0.3333}}@{}}
\toprule\noalign{}
\begin{minipage}[b]{\linewidth}\raggedright
Contribution
\end{minipage} & \begin{minipage}[b]{\linewidth}\raggedright
Value
\end{minipage} & \begin{minipage}[b]{\linewidth}\raggedright
Source
\end{minipage} \\
\midrule\noalign{}
\endhead
\bottomrule\noalign{}
\endlastfoot
Reduced-mass \texttt{πm/M} &
\texttt{+44.66\textquotesingle{}\textquotesingle{}/c} & Corda §§2--4
(\textbf{not a real precession}) \\
Framework relativistic correction \texttt{−πGM/(c²\ r\_0)} &
\textbf{\texttt{−6.87\textquotesingle{}\textquotesingle{}/c}} & wrong
sign vs.~GR's \texttt{+42.99\textquotesingle{}\textquotesingle{}/c} \\
Net \texttt{Δφ\_d} &
\textbf{\texttt{+37.79\textquotesingle{}\textquotesingle{}/c}} & sum,
vs.~observed \texttt{≈\ 43\textquotesingle{}\textquotesingle{}/c} \\
\end{longtable}

The 12\% shortfall (\texttt{+37.79\textquotesingle{}\textquotesingle{}}
vs \texttt{+43\textquotesingle{}\textquotesingle{}}) was interpreted in
PR \#83 as a near-miss. \textbf{Under the present analysis, this is a
coincidence on top of a coincidence}: \texttt{πm/M} is fictitious, and
the relativistic correction has the wrong sign.

\hypertarget{what-pr-83-doc-05-computed-apsidal-angle-the-genuine-answer}{%
\subsubsection{§4.2 What PR \#83 doc 05 computed (apsidal-angle, the
genuine
answer)}\label{what-pr-83-doc-05-computed-apsidal-angle-the-genuine-answer}}

\href{../Mercury_Perihelion/05_nbody_orbital_mechanics.md}{\texttt{Mercury\_Perihelion/05\_nbody\_orbital\_mechanics.md}}
applied Corda's \emph{own} apsidal-angle formula (Eq. 31) to Gill's
modified force law:

\[
F(r) = -\frac{GMm}{r^2}\left(1 - \frac{GM}{c^2 r}\right)
\]

The apsidal-angle formula gives a precession per orbit that is:

\[
\boxed{\;\Delta\varphi_{\rm framework}^{\rm genuine}
       \;=\; -\frac{1}{6}\,\Delta\varphi_{\rm GR}
       \;=\; -7.17\,\text{''/century}\;}
\]

\textbf{Opposite sign, one-sixth magnitude of GR exactly.} This is the
famous ``1/6 factor'' of a force-law-only modification of Newtonian
gravity that lacks GR's spatial-metric curvature. (See
e.g.~Misner-Thorne-Wheeler \emph{Gravitation} §40.6 for the analogous
PPN parameter analysis.)

The N-body refinement cannot close this gap: planetary perturbations
contribute framework relativistic corrections at
\texttt{\textasciitilde{}10\^{}\{-4\}–10\^{}\{-5\}} of the Sun--Mercury
term (since \texttt{GM/(c²r)} scales with the attracting mass and the
other planets are light). The
\texttt{−7.17\textquotesingle{}\textquotesingle{}/century} is the
framework's structural prediction for the eccentric two-body case,
computed via the same apparatus (apsidal angle) Corda himself uses.

\hypertarget{the-recommendation-for-tepper}{%
\subsubsection{§4.3 The recommendation {[}for
Tepper{]}}\label{the-recommendation-for-tepper}}

The Corda-style \texttt{πm/M\ +\ relativistic\ correction} template that
\emph{Relativistic Newtonian Theory} inherits should be replaced by the
apsidal-angle calculation as the framework's genuine perihelion
prediction. The relevant numbers:

\begin{longtable}[]{@{}
  >{\raggedright\arraybackslash}p{(\columnwidth - 4\tabcolsep) * \real{0.3333}}
  >{\raggedright\arraybackslash}p{(\columnwidth - 4\tabcolsep) * \real{0.3333}}
  >{\raggedright\arraybackslash}p{(\columnwidth - 4\tabcolsep) * \real{0.3333}}@{}}
\toprule\noalign{}
\begin{minipage}[b]{\linewidth}\raggedright
Calculation
\end{minipage} & \begin{minipage}[b]{\linewidth}\raggedright
Result
\end{minipage} & \begin{minipage}[b]{\linewidth}\raggedright
Status
\end{minipage} \\
\midrule\noalign{}
\endhead
\bottomrule\noalign{}
\endlastfoot
Corda-style heuristic (PR \#83 doc 03) &
\texttt{+37.79\textquotesingle{}\textquotesingle{}/century} & numerical
near-miss; structurally fictitious \\
Apsidal-angle (PR \#83 doc 05) &
\texttt{−7.17\textquotesingle{}\textquotesingle{}/century} & genuine
precession of modified force law \\
Standard GR &
\texttt{+42.99\textquotesingle{}\textquotesingle{}/century} & observed
\texttt{\textasciitilde{}43\textquotesingle{}\textquotesingle{}/c} \\
\end{longtable}

If the framework's intended Mercury prediction is the apsidal-angle
result, the verdict is: \textbf{wrong sign, \texttt{1/6} magnitude}.
This is the structural signature of any modified-Newtonian gravity that
adds a \texttt{(1\ +\ V/(mc²))}-style factor to the force without
modifying the spatial metric. The 12\% near-miss in PR \#83 doc 03 is an
artefact of taking Corda's \texttt{πm/M} heuristic as the baseline.

\begin{center}\rule{0.5\linewidth}{0.5pt}\end{center}

\hypertarget{specific-questions-for-tepper}{%
\subsection{§5. Specific questions for
Tepper}\label{specific-questions-for-tepper}}

Six questions for your verdict, drawn from the analysis above:

\begin{enumerate}
\def\labelenumi{\arabic{enumi}.}
\item
  \textbf{Bertrand's theorem applicability.} Do you agree that Eq. (31)
  applied to an inverse-square force (including the
  \texttt{(1+m/M)}-rescaled version of Eq. 32) gives
  \texttt{φ\_0\ =\ 2π}, i.e.\textasciitilde no apsidal precession? This
  is the load-bearing critique of §§2--4. If you read Bertrand
  differently or identify a step in Eq. (31)'s derivation that doesn't
  apply, please flag.
\item
  \textbf{Period vs.~precession.} Do you agree that the \texttt{πm/M} in
  Corda's Eqs. (19)--(20) is computing the angle Mercury sweeps in the
  \emph{unperturbed} period \texttt{T\_0} at the \emph{perturbed}
  angular velocity \texttt{ω} --- not the rotation of the orbit's major
  axis between successive perihelia? This is the conceptual distinction
  that determines whether \emph{Relativistic Newtonian Theory} inherits
  a fictitious baseline.
\item
  \textbf{§5 self-refutation.} The same \texttt{πm/M} formula gives
  \texttt{Venus\ 258\textquotesingle{}\textquotesingle{}/c} and
  \texttt{Earth\ 195\textquotesingle{}\textquotesingle{}/c} vs.~observed
  \texttt{8.6\textquotesingle{}\textquotesingle{}/c} and
  \texttt{3.8\textquotesingle{}\textquotesingle{}/c}. Is Corda's
  ``n-body effects'' rationalisation sufficient, or does the empirical
  failure mean the formula is not a perihelion advance? (PR \#83's
  reading is the latter; we'd like your reading recorded.)
\item
  \textbf{§§6--8 vs.~eccentricity.} Corda's \texttt{3π\ r\_g/r\_0}
  reproduces the e=0 GR limit and misses \texttt{1/(1−e²)}. For Mercury
  this is a 4.4\% shortfall. Do you read §§6--8 as a valid pedagogical
  derivation of the circular-orbit GR result (PPN sum of gravitational +
  Doppler TD), or as a complete GR derivation? The distinction matters
  for whether the eccentricity correction needs to be added to the
  framework's prediction.
\item
  \textbf{The apsidal-angle calculation for Gill's force law.} PR \#83
  doc 05's apsidal-angle result
  \texttt{−7.17\textquotesingle{}\textquotesingle{}/century}
  (\texttt{−1/6} of GR exactly) for the framework's force law
  \texttt{a\ =\ −(GM/r²)(1\ −\ GM/(c²r))·ê\_r} is the proper answer to
  ``what does Gill's modified Newton predict for Mercury?''. Do you read
  this as the framework's intended prediction, or is the Corda-style
  heuristic
  (\texttt{+37.79\textquotesingle{}\textquotesingle{}/century})
  intended?
\item
  \textbf{The \texttt{1/6} structural signature.} PR \#83 doc 05
  identifies \texttt{−1/6} of GR as the classic signature of a
  force-law-only modification of Newton without spatial-metric
  curvature. Does this match your understanding of what
  \emph{Relativistic Newtonian Theory} delivers, or does the framework
  include a spatial-metric component that PR \#83 missed?
\end{enumerate}

The answers to these six adjudicate whether \emph{Relativistic Newtonian
Theory}'s Mercury prediction is
\texttt{+37.79\textquotesingle{}\textquotesingle{}/century} (Corda
template), \texttt{−7.17\textquotesingle{}\textquotesingle{}/century}
(apsidal-angle), or something else.

\begin{center}\rule{0.5\linewidth}{0.5pt}\end{center}

\hypertarget{honest-framing-what-this-report-does-and-does-not-claim}{%
\subsection{§6. Honest framing --- what this report does and does not
claim}\label{honest-framing-what-this-report-does-and-does-not-claim}}

\hypertarget{what-this-report-claims}{%
\subsubsection{What this report claims}\label{what-this-report-claims}}

\begin{itemize}
\tightlist
\item
  Corda's algebra is mostly correct (specifically: Eqs. 17, 20, 31, 37,
  67, 77 are all Wolfram-MCP verified).
\item
  Corda's §§2--4 result \texttt{Δφ\ =\ πm/M} is \emph{not} a perihelion
  advance --- it is a period-and-angular-velocity rescaling that does
  not change the apsidal angle of a closed inverse-square orbit. This is
  verified via Corda's own Eq. (31).
\item
  The §5 numerical breakdown for Venus and Earth is empirical evidence
  that the formula is not a real perihelion advance.
\item
  §§6--8 reproduce the GR result in the circular-orbit limit only.
\item
  \emph{Relativistic Newtonian Theory}'s
  \texttt{+37.79\textquotesingle{}\textquotesingle{}/century} prediction
  inherits the Corda framing's fictitious \texttt{πm/M} baseline plus a
  wrong-sign relativistic correction.
\item
  The framework's \emph{genuine} perihelion prediction, computed via the
  apsidal-angle formula on the modified force law, is
  \texttt{−7.17\textquotesingle{}\textquotesingle{}/century}
  (\texttt{−1/6} of GR).
\end{itemize}

\hypertarget{what-this-report-does-not-claim}{%
\subsubsection{What this report does not
claim}\label{what-this-report-does-not-claim}}

\begin{itemize}
\tightlist
\item
  That Corda's paper is not a valid published work. It is a
  peer-reviewed paper in \emph{Physics of the Dark Universe} (a
  respectable Elsevier journal). The findings here are about the
  \emph{interpretation} of correctly-derived algebra.
\item
  That \emph{Relativistic Newtonian Theory} as a whole is wrong. The
  Mercury perihelion is one test; other tests (the four candidate
  curved-spacetime extensions tracked in the GPS author report's Q1)
  remain open.
\item
  That the framework cannot deliver a Mercury-matching prediction. The
  framework may have alternative extensions to gravity beyond the one
  analysed in \emph{Relativistic Newtonian Theory}; this report only
  addresses the specific extension that paper proposes.
\item
  That the apsidal-angle critique is novel. It is straightforward
  textbook orbital mechanics. The novelty here is applying it explicitly
  to Corda's framework and to \emph{Relativistic Newtonian Theory} as an
  inheritance check.
\end{itemize}

\hypertarget{what-the-report-is-asking-from-you}{%
\subsubsection{What the report is asking from
you}\label{what-the-report-is-asking-from-you}}

Six adjudications (§5 above). Your verdict on these determines what the
campaign's Mercury verdict should read; the campaign's current PR \#83
documents both the Corda-style heuristic and the apsidal-angle critique
and is held pending your reading.

\begin{center}\rule{0.5\linewidth}{0.5pt}\end{center}

\hypertarget{cross-references}{%
\subsection{§7. Cross-references}\label{cross-references}}

\begin{itemize}
\tightlist
\item
  \textbf{This document}:
  \href{2026-05_corda_perihelion_review_for_gill.md}{\texttt{Roadmapping/Author\_Reports/2026-05\_corda\_perihelion\_review\_for\_gill.md}}
\item
  \textbf{Corda PDF}:
  \href{../External_Papers/Secret_Corda.pdf}{\texttt{Roadmapping/External\_Papers/Secret\_Corda.pdf}}
\item
  \textbf{Corda markdown} (marker-pdf):
  \href{../Converted_Markdown/Secret_Corda/Secret_Corda.md}{\texttt{Roadmapping/Converted\_Markdown/Secret\_Corda/Secret\_Corda.md}}
\item
  \textbf{Concise equation-verification doc}:
  \href{../Equation_Verification/Secret_Corda_Perihelion.md}{\texttt{Roadmapping/Equation\_Verification/Secret\_Corda\_Perihelion.md}}
\item
  \textbf{Gill's \emph{Relativistic Newtonian Theory} analysis (PR
  \#83)}:

  \begin{itemize}
  \tightlist
  \item
    PR: https://github.com/temoTxt/PyPhysics/pull/83
  \item
    \href{../Equation_Verification/Dual_Newtonian_Theory.md}{\texttt{Roadmapping/Equation\_Verification/Dual\_Newtonian\_Theory.md}}
  \item
    \href{../Mercury_Perihelion/README.md}{\texttt{Roadmapping/Mercury\_Perihelion/README.md}}
  \item
    \href{../Mercury_Perihelion/05_nbody_orbital_mechanics.md}{\texttt{Roadmapping/Mercury\_Perihelion/05\_nbody\_orbital\_mechanics.md}}
    (apsidal-angle calculation; the \texttt{−1/6} of GR result)
  \end{itemize}
\item
  \textbf{GPS author report Q1 (curved-spacetime extension question)}:
  \href{2026-05_gps_relativity_summary_for_gill.md}{\texttt{Roadmapping/Author\_Reports/2026-05\_gps\_relativity\_summary\_for\_gill.md}}
\item
  \textbf{Issue \#81}: https://github.com/temoTxt/PyPhysics/issues/81
\end{itemize}

\hypertarget{human-acceptance-section-crocco-5-substantive-ai}{%
\subsection{§8. Human-acceptance section (Crocco §5 substantive
AI)}\label{human-acceptance-section-crocco-5-substantive-ai}}

\end{document}
